\chapter{Experimental Native C API}
\label{ch:native-c-api}
\index{C API}\index{ABI version}\index{CNA::CApi}

\CnaEditionTag{} contains the source and public headers for an experimental C17 binding over CNA's
canonical C\texttt{++} implementation. It is intended as an optional shared-library surface, not
a replacement implementation and not evidence that every C\texttt{++} expression has a useful C
spelling. The tag-level build result is important context for the nominal command:

\begin{lstlisting}[style=cnashell]
cmake -S . -B build-c \
  -DCNA_BUILD_C_API=ON \
  -DCNA_GRAPHICS_RENDERER=HEADLESS \
  -DCNA_PLATFORM=HEADLESS \
  -DCNA_AUDIO_PLATFORM=NULL \
  -DCNA_ENABLE_NET=ON
cmake --build build-c --target cna_c_api
\end{lstlisting}

\begin{limitationnote}
This command does not produce the library at the immutable alpha.1 tag. A clean GNU~14.2 build
with networking enabled reaches \path{CnaCApiCoreExt.cpp} and fails its own compile-time identity
invariant: the C table and \texttt{CNA\_GRAPHICS\_RENDERER\_MAXIMUM} contain 49 identities, while
the canonical C\texttt{++} enumeration contains 50. The missing C identity is \texttt{NANOVG}.
This is a build-blocked source surface, not a consumable alpha.1 binary artifact. The command is
retained so the failure is reproducible and a later tag can be rechecked against the same build
contract.
\end{limitationnote}

The target definition is \texttt{cna\_c\_api}, exported to CMake consumers as
\texttt{CNA::CApi}. Its intended normal product is a shared library; the source also defines an
explicitly controlled static form for packaging tests. Public headers live under \path{CNA/C/}, and
\path{CNA/C/cna.h} is the 59-header umbrella. A C consumer includes those headers and links the C
API library; it does not include CNA or Sharp Runtime C\texttt{++} headers.
The project builds its C translation units as required C17.  The compatibility generator also
checks that each public header remains self-contained under C99, C11, C17 and C23 modes where the
declared compiler is installed; C99 header compatibility does not lower the project's C17 build
standard.

\section{Two version systems, deliberately independent}

\begin{center}
\small
\begin{tabularx}{\textwidth}{>{\raggedright\arraybackslash}p{3.2cm}p{3.4cm}X}
\toprule
\textbf{Version} & \textbf{At this tag} & \textbf{What it governs} \\
\midrule
CNA product & \CnaProductVersion{} & Source release, C\texttt{++} API, tools, modules and
  documentation. The leading zero and prerelease component mean that compatibility may still
  change. \\
Native C ABI & \texttt{0.7.0} & Binary layouts, exported C symbols and the C binding's
  compatibility rules. \texttt{CNA\_ABI\_VERSION} encodes it as a 32-bit integer. \\
\bottomrule
\end{tabularx}
\end{center}

The header macros \texttt{CNA\_ABI\_VERSION\_MAJOR}, \texttt{\_MINOR}, \texttt{\_PATCH} and
\texttt{CNA\_ABI\_VERSION} describe the headers. \texttt{cna\_get\_abi\_version()} reports the
loaded library. Consumers should compare the encoded values and reject an incompatible major.
This is separate from the generated C\texttt{++} product-version API in
\path{CNA/Version.hpp}; \texttt{0.7.0} is not CNA product release 0.7.0, and
\CnaProductVersion{} is not C ABI 0.1.0.

\begin{limitationnote}
The tag's generated C API release-gate report still labels its header ``CNA C ABI 0.1.0'', while
\path{CNA/C/abi.h} and \texttt{cna\_get\_abi\_version()} establish 0.7.0. The implementation and
public header are the edition authority; the report label is a documentation defect, not a second
supported ABI. Its ``Ready'' verdict is not a build result: the checker confirms that test,
install and consumer-gate definitions exist, but does not execute them, and therefore misses the
49-versus-50 compile failure.
\end{limitationnote}

\section{Results, diagnostics, and handles}

Every fallible operation returns the fixed-width \texttt{CNA\_Result}. At the tag the public
codes are 0 through 14: success, invalid argument, invalid handle, invalid state, out of memory,
I/O, unsupported, platform, thread, callback, overflow, encoding, internal, shutting down, and
buffer too small. Failure details are thread-local. Call
\texttt{cna\_error\_get\_last\_info()}, query the UTF-8 message size, then copy the message into
caller-owned storage. A successful call can replace the diagnostic, so error retrieval belongs
immediately beside the failing call.

Objects cross the boundary as the 64-bit \texttt{CNA\_Handle}; zero is invalid. The internal
registry validates generation, object kind, owning runtime, lifetime and thread affinity before
recovering a C\texttt{++} object. Destroy operations invalidate their handle. A borrowed handle
does not transfer ownership, callbacks must obey the route's re-entry rule, and a handle is not a
native pointer that C may dereference. Versioned structures carry \texttt{struct\_size} and
\texttt{struct\_version}; strings use explicit UTF-8 byte lengths; collections use counts plus
indexed access or dedicated enumerator handles.

\section{Surface and coverage boundary}

The headers cover framework/runtime, math and value types, graphics devices and resources,
effects and models, input, content, gamer services and networking, storage, audio/XACT, media,
sensors, devices, components, windows, and the graphics-device manager. This breadth should not
be summarized as ``complete C coverage.'' The tag's generated inventory classifies 6,712 public
C\texttt{++} declarations as follows:

\begin{center}
\small
\begin{tabular}{lr}
\toprule
\textbf{Disposition} & \textbf{Declarations} \\
\midrule
Implemented & 6,317 \\
Partially mapped, with an approved limitation & 15 \\
Planned / unclassified & 0 \\
No C form & 380 \\
\midrule
Total & 6,712 \\
\bottomrule
\end{tabular}
\end{center}

The 380 exclusions include deleted operations, iterators, arbitrary C\texttt{++} templates,
protected subclass hooks, Sharp Runtime objects, friendship and internal platform interfaces.
The 15 partial entries name their usable subset and alternative route. For example, generic
\cnaclass{ContentManager::Load} cannot accept an arbitrary C\texttt{++} type from C, so the binding
supplies typed asset routes. This is reviewed coverage, not one-to-one syntactic transliteration.

The generated limitations prose opens with a stale total of 6,415 even though its own current
rows and the coverage inventory sum to 6,712. This edition uses the reproducible inventory and
records the prose total as drift.

\section{Renderer and platform consequences}

Once the adapter builds, its C routes create canonical CNA objects, so their capabilities and
lifetime rules remain those of the active renderer and selected platform/audio modules. It does not expose an
\cnaclass{IGraphicsRenderer} pointer or a second renderer implementation. A multi-renderer build
still resolves the runtime renderer through CNA's C\texttt{++} selection service before the first
graphics device; C routes then observe the chosen device and receive
\texttt{CNA\_RESULT\_NOT\_SUPPORTED} for unavailable operations. Renderer rebuild for a changed
multisample request is explicitly resource-destructive and is only safe before GPU resources
exist.

This corrects an ambiguity in the tag's C API renderer prose, which says only that the backend is
``compiled into CNA.'' That wording describes a single-renderer build but does not erase the
tag's generated multi-renderer registry.

\section{Verification inventory, not a universal pass claim}

Source inspection at the tag finds 74 pure-C test translation units, 82 CTest registrations whose
names begin \texttt{CApi\_}, 176 recorded structure layouts, and 2,861 recorded exports. Five
dedicated workflows cover the ABI baseline, compatibility matrix, coverage inventory,
limitations matrix, and combined release gate. Their presence proves automation is declared;
it does not by itself prove that GitHub executed every cell for this tag. The compatibility
matrix describes 23 compiler cells across seven toolchain families, while actual outcomes require
workflow-run evidence.

The declared release gates are designed to check that public declarations have a disposition,
generated reports are current, headers compile as C and C\texttt{++}, layout/export baselines are
reviewable, examples consume the installed package, and selected renderer builds validate the
same ABI boundary. At alpha.1, neither the five narrow workflows nor the source-only aggregate
checker closes the final-library-build predicate. The missing NanoVG identity, stale ABI label and
stale declaration count show why generated records still require a fresh artifact build.

\begin{evidencenote}
The header count, result values, ABI 0.7.0, coverage partition, test-source and registration counts
are Source-proven at \CnaEditionTag{}. They are not reported as an executed 82-test run. Two
build probes were attempted. With \texttt{CNA\_ENABLE\_NET=OFF}, configuration succeeds but the C
adapter cannot see GamerServices headers because CMake neither refuses nor repairs the invalid
closure. With networking enabled, compilation reaches and fails the 49-versus-50 renderer
\texttt{static\_assert}. The tag-level library status is therefore \textbf{Blocked at compile
time}, while the generated product-version example remains independently Compile-proven.
\end{evidencenote}
